\documentclass[10pt]{article}
\usepackage[a4paper,margin=22mm]{geometry}
\usepackage[T1]{fontenc}
\usepackage{lmodern,amsmath,amssymb,amsthm,mathtools,microtype,graphicx}
\usepackage[hidelinks]{hyperref}
\newtheorem{theorem}{Theorem}
\newtheorem{corollary}[theorem]{Corollary}
\title{A complete fixed-target cofactor classification\ and original Erd\H{o}s--Straus channel returns}
\author{The Clankers}
\date{20 September 2026}
\begin{document}
\maketitle
\begin{abstract}
We classify all divisors of $j^2$ in a fixed positive class $t$ modulo $4j-1>t$.
A canonical divisor and an unbounded companion are separated from a finite
exceptional set with sharp grade bound $4j-1\le t^2-3t+1$. The bound is realized
by original exterior states at infinitely many hard primes. Applied to original
negative shape coordinates, the classification gives complete same-grade middle
returns and forced availability on entire nonsquare shape lines. It also proves
that each fixed alpha coefficient $t\ge4$ has infinitely many original hard-prime
examples with no same-grade middle return. This is not a universal ES proof.
\end{abstract}

\begin{theorem}[Complete capacity and sharp boundary]
For $h=4j-1>t\ge1$, define
$\mathcal A_t(j)=\{W\mid j^2:W\equiv t\pmod h\}$.
Its disjoint, uniquely labelled words are
\[
 W=t\quad(t\mid j^2),\qquad
 W=4tj\quad(4t\mid j),                                  \tag{1}
\]
and the finite words
\[
 1\le\ell\le\lfloor(t-2)/4\rfloor,\quad
 w\mid\ell^2,\quad n=(t-w)/(4\ell+1)\in\mathbb Z_{>0},
\]
\[
 V=\ell^2/w,\quad j=\ell+4nV,\quad h>t,
 \qquad W=j^2/V=t+nh.                                  \tag{2}
\]
Every finite word satisfies $h\le t^2-3t+1$. Equality holds exactly at
$w=n=1$, $t=4\ell+2$. In particular, a nonempty box with $t\nmid j^2$ is in
this finite grade range. Beyond it the exact count is
$\mathbf1_{t\mid j^2}+\mathbf1_{4t\mid j}$.
\end{theorem}
\begin{figure}[htbp]
\centering
\includegraphics[width=\textwidth]{figures/fixed-target-capacity-partition.pdf}
\caption{Exact inverse classification in Theorem 1.  The canonical,
companion and finite branches are disjoint, and every label is recovered from
the original divisor \(W\).}
\end{figure}
\begin{proof}
For $W\ne t$, write $W=t+nh$, $n\ge1$, and $V=j^2/W$. The identity
$16WV=(h+1)^2$ gives $16tV=1+kh$, where $k\ge1$ and $k\equiv1\pmod4$.
Put $k=4\ell+1$. Subtraction gives
\[
 j=\ell+4nV,\qquad t=(4\ell+1)n+\ell^2/V.
\]
At $\ell=0$, $n=t$ and $j=4tV$, exactly the second word of (1). At $\ell>0$,
$w=\ell^2/V$ is a positive integer and $t\ge4\ell+2$. All steps reverse,
with inverse $n=(W-t)/h$, $V=j^2/W$, $\ell=j-4nV$. Thus the list has no
omissions or duplicate labels. Finally
\[
 h+1=4\ell+\frac{16\ell^2(t-w)}{(4\ell+1)w}
 \le4\ell+\frac{16\ell^2(t-1)}{4\ell+1}
 \le t^2-3t+2.
\]
The last expression increases strictly in $\ell$, whose maximum is $(t-2)/4$.
Equality conditions give $w=n=1$. The companion always coexists with the
canonical word, so it cannot repair its absence.
\end{proof}

The sharp family is
\[
 t=4\ell+2,\quad j=\ell(4\ell+1),\quad
 h=t^2-3t+1,\quad W=(4\ell+1)^2,\quad V=\ell^2.           \tag{3}
\]
Here $t\nmid j^2$, because $2\ell+1>1$ is coprime to $j$. Finite words can
also coexist with the canonical one: $\mathcal A_{289}(1224)=\{289,5184\}$.
The companion in (1) is genuinely unbounded and must not be folded into the
finite exceptional claim.

\paragraph{Original arithmetic and inverse.}
For prime $p\equiv1\pmod4$, an original exterior state has
\[
 p/4<a<p/2,\quad R=4a-p,\quad u\mid a^2,\quad R\mid4u+1.
\]
Normalize by $g=\gcd(a,u)$, $(h,r,s)=(g^2/u,u/g,a/g)$ and
$\kappa=(pr+s)/R$. The original ordered return is
$(a,hs\kappa,phr\kappa)$; its inverse is $u=a^2/(Ry-pa)$.
Keep
\[
 v=hs^2,\quad D=(4u+1)/R,\quad
 \alpha=v-R,\quad\beta=v-D.
\]
The original identities yield
$h\mid\alpha\beta-1$, $p\equiv\alpha\pmod h$, $p\beta\equiv1\pmod h$.
Reciprocity on the actual factors gives $(h/p)=-1$.

If $\alpha=-4t$ or $\beta=-4t$, then
$h\equiv3\pmod4$, $s$ is odd, $\gcd(h,t)=1$ and $(t/h)=1$.
Set $j=(h+1)/4$. The entire middle cofactor-$h$ gate is exactly
\[
 W\in\mathcal A_t(j),\quad
 U=W\ (\alpha),\qquad U=j^2/W\ (\beta).
\]
For every word it returns
\[
 R_M=(p+4U)/h,\quad a_M=(p+R_M)/4,\quad g_M=\gcd(j,U),
\]
\[
 h_M=g_M^2/U,\quad r_M=U/g_M,\quad\lambda_M=j/g_M,
 \quad s_M=R_M\lambda_M-r_M,                             \tag{4}
\]
with denominators $(a_M,ph_Ms_M\lambda_M,ph_Mr_M\lambda_M)$.
Indeed $U\mid j^2$, $h\mid p+4U$, and
$4U\le(h+1)^2/4<(h-1)p$ since $p>2h$. Thus $0<R_M<p$ and
$R_M\equiv3\pmod4$. Prime valuations give the integral normalization;
$a_M=R_Mj-U=h_Mr_Ms_M>0$, and a common divisor of $r_M,R_M$ would divide
$p$, although $r_M\le j<p$. This proves the original coprimality and both range
bounds. Complementing $U$ exchanges the retained middle orientation.

These are classical Type-II coordinates \cite[\S2]{ET}; the contribution is the
complete capacity input and its original source supply. In (2), putting
$g_0=\gcd(\ell,V)$ gives $H_0=g_0^2/V$, $r_0=V/g_0$, $\lambda_0=\ell/g_0$.
The returned $(h_M,r_M,\lambda_M)$ equals
$(H_0,\lambda_0+4nr_0,r_0)$ or $(H_0,r_0,\lambda_0+4nr_0)$.
Thus $h_M\le\ell<t/4$ on the finite branch.

\begin{theorem}[Automatic entire shape lines]
At a prime $p\equiv1\pmod{24}$, the original source forces $t\mid j^2$ on
\[
 \alpha=-4t,\quad t\in\{1,2,3\};\qquad
 \beta=-4t,\quad t\mid36.
\]
Every such source therefore has the explicit middle return (4), without any
bound on its other shape coordinate.
\end{theorem}
\begin{proof}
For alpha, $(2/h)=1$ with $h\equiv3\pmod4$ forces $h\equiv7\pmod8$;
$(3/h)=1$ forces $h\equiv11\pmod{12}$. These give $2\mid j$ or $3\mid j$.
For beta, $D=hs^2+4t$, $r+\kappa=sD$ odd and $pD+1=4hr\kappa$ give
$D\equiv7\pmod8$ and $h\equiv7-4t\pmod8$. If $2\mid t$, then $2\mid j$.
If $3\mid t$, all equalities in the remainder of this sentence are modulo
three: the case $h=0$ is impossible; at $h=1$, $s=0$ would
require $1=-r^2$; $s=\pm1$ would require $2=r(s-r)$, also impossible.
Hence $h=2\pmod3$ and $3\mid j$. Since each exponent in $t\mid36$ is at
most two, this proves the full divisor budget.
\end{proof}

\paragraph{Sharpness in the original prime domain.}
For every $t=6+420z$, put $h=t^2-3t+1$, $R=t^2+t+1$.
Then $\gcd(hR,210)=\gcd(h,R)=1$. The integer $r_0=(t+2R)/4$ satisfies
$16tr_0^2\equiv1\pmod R$, since $t^3\equiv1\pmod R$.
Solve
$r\equiv r_0\pmod R$ and
$r\equiv((R+1)/4)h^{-1}\pmod{210}$.
The progression $p=4hr-R$ is reduced modulo $840hR$ and is 1 modulo 840.
Dirichlet's theorem \cite[Theorem 18.1]{Dirichlet} therefore gives infinitely
many prime members. They have the original E state $(a,u)=(hr,hr^2)$,
$\alpha=-4t$, and the unavailable canonical word but available sharp word (3).
For example the hard prime $825241$ has
\[
 (a,R,u,h,r,s,\kappa)=(206321,43,2240439739,19,10859,1,208402140).
\]
At $j=5$, only $W=25$ in $\{1,5,25\}$ has residue six modulo 19. It returns
\[
 (a_M,R_M,U,h_M,r_M,s_M,\lambda_M)
 =(217170,43439,25,1,5,43434,1).
\]
\begin{figure}[htbp]
\centering
\includegraphics[width=\textwidth]{figures/capacity-sharp-bound.pdf}
\caption{The sharp finite boundary and the exact \(p=825241\) nonsquare
return.  Blue points enumerate every proved finite template through \(t=120\);
the displayed quadratic curve is the theorem.}
\end{figure}

\paragraph{A complete limitation of same-grade repair.}
For every fixed $t\ge4$ there are infinitely many actual hard-prime E states
with $\alpha=-4t$ and $\mathcal A_t(j)=\varnothing$.
Here is the constructive proof. Put
$K=\prod_{l^e\parallel t}l^{\lceil e/2\rceil}$ over every prime, including
two, and \(M=4K\). For odd \(A\) coprime to \(t\), let
\(\chi_t(A)=(t/A)\) be the Kronecker--Jacobi character determined by the
squarefree part of \(t\). The unit classes \(A\equiv3\pmod4\),
\(\chi_t(A)=1\) number \(\varphi(M)/2\)
if $t$ is a square and $\varphi(M)/4$ otherwise. The latter count follows by
CRT: the positive quadratic character of $t$ is nontrivial on units 1 modulo
four and has value one at $-1$. This class set is the singleton $\{-1\}$ only
at $t=1,2,3$. Choose $A\ne-1$ and lift it so that $R,R-4t$ are units at
$3,5,7$. The added restrictions each exclude at most two residues, and are
already satisfied when that prime divides $K$.
Dirichlet supplies a prime \(R\) in this reduced class with
$h=R-4t>\max(4t,t^2-3t+1)$. Choose \(q^2\equiv t\pmod R\), put
\(r_0\equiv(4q)^{-1}\pmod R\), and impose the
same CRT on $r$ as above. Again $p=4hr-R$ is a reduced progression of primes
$1\pmod{840}$. Its original E state is $(hr,hr^2)$.
But $h\not\equiv-1\pmod{4K}$ makes $t\nmid j^2$, and Theorem 1 excludes
all other words at this grade. Every such prime already has its constructed
E solution; this is not an ES counterexample.

For $t=4$, take $R=59,h=43,j=11$ and
$p=1447321+2131080n$. The prime member $n=3$ is $7840561$.
The entire capacity set tests $1,11,121$ modulo 43, giving $1,11,35$ rather
than four. The beta automatic list cannot be extended indiscriminately:
actual hard-prime beta sources with coefficients $5,8,27$ also have empty
old-grade boxes; full certificates are included.

\paragraph{Scope and exact records.}
The full note gives every inverse-source fibre, including its original exterior
gate, and corrects a missing gate in the predecessor's prose (its executable
already tested it). It preserves every word and ordered denominator. The
standard-library programs compare literal divisor boxes and independently
constructed primitive source states. Their declared bounds and execution
receipts distinguish finite replay from the unbounded proofs above. No initial
ES state has been forced at an arbitrary prime; universal occupancy, historical
priority, a Lean build, and independent mathematical review are not claimed.

\begin{thebibliography}{9}
\bibitem{ET} Elsholtz, C., \& Tao, T. (2013). Counting the number of solutions to
the Erd\H{o}s--Straus equation on unit fractions. \emph{J. Aust. Math. Soc.},
94, 50--105. \texttt{arXiv:1107.1010}, \S2.
\bibitem{Dirichlet} Sutherland, A. V. (2017). \emph{18.785 Number Theory I,
Lecture 18}, MIT, Theorem 18.1.
\end{thebibliography}
\end{document}
